\documentclass[11pt,a4paper]{article}

\usepackage[utf8]{inputenc}
\usepackage[T1]{fontenc}
\usepackage[spanish,es-tabla]{babel}
\usepackage[a4paper,margin=2.5cm]{geometry}
\usepackage{booktabs}
\usepackage{tabularx}
\usepackage{array}
\usepackage{titlesec}
\usepackage{hyperref}
\usepackage{graphicx}

\hypersetup{
  colorlinks=true,
  linkcolor=black,
  urlcolor=blue!50!black,
  citecolor=blue!50!black
}

\titleformat{\section}{\normalfont\Large\bfseries}{\thesection.}{0.6em}{}
\titleformat{\subsection}{\normalfont\large\bfseries}{\thesubsection}{0.6em}{}

\begin{document}

\begin{titlepage}
\begin{center}

\vspace*{1cm}

{\large UNIVERSIDAD TÉCNICA ESTATAL DE QUEVEDO}\\[0.2cm]
{\normalsize Facultad de Ciencias de la Computación}\\
{\normalsize Carrera de Ingeniería de Software}\\
{\normalsize Aplicaciones Distribuidas (ISR-701)}

\vspace{2cm}

\rule{\linewidth}{0.4pt}\\[0.4cm]
{\huge \bfseries Informe de Revisión Técnica Cruzada}\\[0.3cm]
{\large \bfseries del Código del Proyecto Fin de Curso}\\[0.4cm]
\rule{\linewidth}{0.4pt}

\vspace{0.8cm}
{\normalsize Taller FORM-TL-05 — Unidad 5\\
Principios SOLID, patrones de diseño, separación de capas,\\
manejo de excepciones distribuidas y observabilidad}

\vspace{2cm}

\begin{tabular}{r l}
\textbf{Equipo revisor:} & BCEL \\
\textbf{Sistema revisado:} & SCLI — Sistema de Control de Laboratorios e Infraestructura \\
\textbf{Equipo autor del PFC:} & FUVV \\
\textbf{Fecha de auditoría:} & 11 de agosto de 2026 \\
\end{tabular}

\vfill

\textbf{Integrantes del equipo revisor:}\\[0.3cm]
\begin{tabular}{l l}
\textbf{Castro López Pedro Leonardo} & Líder / Editor / Revisión Bloques C, D, E \\
\textbf{Emanuel Pino Juliana Romina} & Revisión Bloques A y B (SOLID y Patrones) \\
\end{tabular}

\vspace{1cm}
{\small Quevedo — Los Ríos — Ecuador}

\end{center}
\end{titlepage}
\section{Objeto auditado}

La revisión se realizó sobre un commit específico del repositorio del equipo FUVV,
fijado antes de cualquier análisis para garantizar la reproducibilidad de los
hallazgos.

\begin{center}
\begin{tabular}{l l}
\toprule
\textbf{Repositorio revisado} & \url{https://github.com/iavillamarin98-pred/Proyecto_Fin_de_Curso_scli} \\
\textbf{Rama} & main \\
\textbf{Commit auditado (SHA completo)} & \texttt{49ec17dfd42835fe5dfbd75249a5e8a51941ef4a} \\
\textbf{SHA corto} & \texttt{49ec17d} \\
\textbf{Fecha del commit} & 2026-07-29 23:40:45 (UTC-05) \\
\textbf{Autor del commit} & Freddy Vladimir Farinango Guandinango \\
\textbf{Archivos versionados} & 414 \\
\textbf{Fecha y hora de la auditoría} & 2026-08-11 \\
\bottomrule
\end{tabular}
\end{center}

\section{Fase 1: Congelación del objeto revisado}

Antes de iniciar cualquier análisis se ejecutaron los comandos exigidos por la
guía del taller para dejar registrada la versión exacta del código sobre la que
se produce toda la evidencia posterior. La salida obtenida fue la siguiente:

\begin{verbatim}
$ git log -1 --format="%H | %aI | %an"
49ec17dfd42835fe5dfbd75249a5e8a51941ef4a | 2026-07-29T23:40:45-05:00 |
Freddy Vladimir Farinango Guandinango

$ git rev-parse --short HEAD
49ec17d

$ git ls-files | wc -l
414
\end{verbatim}

El sistema levanta correctamente mediante \texttt{docker compose up -d}
conforme al \texttt{README.md} del repositorio auditado. No se detectaron
desviaciones que impidieran la revisión.
